\chapter{โครงงานบูรณาการและการพัฒนาต้นแบบนวัตกร AIoT ภาคสนาม}
\label{chap:ch10}

\section*{แผนบริหารการสอนประจำบทที่ 10}
\addcontentsline{toc}{section}{แผนบริหารการสอนประจำบทที่ 10}

\begin{planbox}{หัวข้อเนื้อหาประจำบท}
\begin{enumerate}[topsep=2pt, itemsep=1pt, partopsep=0pt, parsep=0pt, leftmargin=1.5em]
    \item สถาปัตยกรรมไปป์ไลน์บูรณาการระบบนวัตกร AIoT สวนทุเรียนพรีเมียมครบวงจร
    \item สถาปัตยกรรมไปป์ไลน์บูรณาการระบบนวัตกร AIoT ฟาร์มสัตว์น้ำอัจฉริยะครบวงจร
    \item ระเบียบวิธีตรวจสอบความใช้ได้ของวิธีทางสถิติและการทดสอบภาคสนามในพื้นที่จริง
    \item ผลการประหยัดพลังงานและการเพิ่มมูลค่าผลผลิตในแปลงทดสอบจังหวัดจันทบุรีและตราด
    \item รูบริกเกณฑ์การประเมินผลโครงงานบูรณาการและทิศทางการวิจัยต่อยอดในอนาคต
\end{enumerate}
\end{planbox}

\begin{planbox}{วัตถุประสงค์เชิงพฤติกรรม}
\begin{enumerate}[topsep=2pt, itemsep=1pt, partopsep=0pt, parsep=0pt, leftmargin=1.5em]
    \item บูรณาการระบบฮาร์ดแวร์ เซนเซอร์ ปัญญาประดิษฐ์ และคลาวด์เข้าเป็นระบบเดียวกันได้
    \item ดำเนินการทดสอบความถูกต้องและเสถียรภาพของระบบในแปลงเกษตรจริงได้อย่างสมบูรณ์
    \item วิเคราะห์ผลการทดลองทางสถิติด้วยการทดสอบทีและแผนภาพบลานด์-อัลต์แมนได้
    \item สรุปผลสัมฤทธิ์และนำเสนอโครงงานตามเกณฑ์มาตรฐานวิชาการได้อย่างมืออาชีพ
\end{enumerate}
\end{planbox}

\newpage

\begin{planbox}{วิธีสอนและกิจกรรมการเรียนการสอน}
\begin{enumerate}[topsep=2pt, itemsep=1pt, partopsep=0pt, parsep=0pt, leftmargin=1.5em]
    \item บรรยายเชิงระบบการเชื่อมต่อสถาปัตยกรรมต้นแบบสมบูรณ์ (End-to-End Pipeline)
    \item แบ่งกลุ่มผู้เรียนปฏิบัติการทดสอบระบบบูรณาการในสถานีจำลองภาคสนาม
    \item การวิเคราะห์ข้อมูลทางสถิติและประเมินความคลาดเคลื่อนเทียบกับวิธีมาตรฐาน
    \item การนำเสนอโครงงานบูรณาการ (Capstone Project Presentation) และการสาธิตหน้างาน
\end{enumerate}
\end{planbox}

\begin{planbox}{สื่อการเรียนการสอน}
\begin{enumerate}[topsep=2pt, itemsep=1pt, partopsep=0pt, parsep=0pt, leftmargin=1.5em]
    \item เอกสารคำสอนบทที่ 10 สไลด์บรรยาย และคู่มือเกณฑ์การประเมินรูบริกโครงงาน
    \item ชุดสถานีทดสอบระบบบูรณาการสวนทุเรียนและฟาร์มสัตว์น้ำจำลอง
    \item เครื่องมือวิเคราะห์ทางสถิติและซอฟต์แวร์ประมวลผลข้อมูล (Python SciPy/Statsmodels)
    \item แดชบอร์ดสรุปผลการทำงานของระบบแบบเรียลไทม์
\end{enumerate}
\end{planbox}

\begin{planbox}{การวัดและประเมินผล}
\begin{enumerate}[topsep=2pt, itemsep=1pt, partopsep=0pt, parsep=0pt, leftmargin=1.5em]
    \item การประเมินผลงานโครงงานบูรณาการตามเกณฑ์รูบริกสัดส่วนร้อยละ 50
    \item การประเมินทักษะการแก้ปัญหาทางวิศวกรรมเฉพาะหน้าในแปลงทดสอบภาคสนาม
    \item การประเมินคุณภาพเอกสารรายงานโครงงานฉบับสมบูรณ์และการนำเสนอต่อสาธารณะ
\end{enumerate}
\end{planbox}

\clearpage

\section{สถาปัตยกรรมไปป์ไลน์บูรณาการระดับปฏิบัติการ}

การพัฒนานวัตกร AIoT ไม่ใช่เพียงการต่อวงจรอุปกรณ์แยกส่วน แต่เป็นการผสานเทคโนโลยีหลายระดับเข้าเป็นระบบอัจฉริยะอัตโนมัติครบวงจร (End-to-End Autonomous System)

\begin{table}[H]
\centering
\caption{สถาปัตยกรรมไปป์ไลน์บูรณาการของทั้ง 2 กิจกรรมหลัก}
\label{tab:pipeline_matrix}
\begin{tabularx}{\textwidth}{p{2.8cm}YY}
    \toprule
    \textbf{ชั้นสถาปัตยกรรม} & \textbf{กิจกรรมที่ 1: สวนทุเรียนพรีเมียม} & \textbf{กิจกรรมที่ 2: ฟาร์มสัตว์น้ำอัจฉริยะ} \\
    \midrule
    1. ตรวจวัดกายภาพ & เซนเซอร์ดิน 7-in-1, VPD อากาศ, Tensiometer & หัววัด Optical DO, ความเค็ม, อุณหภูมิ, pH \\
    2. ควบคุมฮาร์ดแวร์ & โซลินอยด์วาล์ว 24VAC หน่วงปิดกันค้อนน้ำ & อินเวอร์เตอร์ VFD ขับมอเตอร์ 3 เฟส เครื่องตีน้ำ \\
    3. ชลศาสตร์/อุณหพลฯ & คำนวณ $ET_c$ ฟีโนโลยี และสมการ Hazen-Williams & คำนวณ $C^*_s$ Benson-Krause และอัตรา OTR \\
    4. โครงข่ายสื่อสาร & ESP-NOW ในสวน เกตเวย์ MQTT LINE แจ้งเตือน & LoRaWAN AS923 ข้ามผืนน้ำ ทนไอเค็มทะเล \\
    5. คอมพิวเตอร์วิทัศน์ & ตัดสีหนามร่องพู HSV และสเปกตรัมเคาะเสียง & วิทัศน์ยอยกอาหารกุ้ง คำนวณอัตราส่วนเม็ดอาหาร \\
    6. ปัญญาประดิษฐ์เอดจ์ & YOLOv8 คัดเกรดผลไม้และตรวจโรคผลเน่า & MobileNetV3 จำแนกโรคกุ้ง 4 อาการ และสเปกตรัม \\
    7. การขับเคลื่อนอัตโนมัติ & ก้านสูบลมดีดผลตกเกรดบนสายพานคัดแยก & ออโต้ฟีดเดอร์วงปิด และสลับมอเตอร์ Lead-Lag 4 สเต็ป \\
    \bottomrule
\end{tabularx}
\end{table}

\section{การตรวจสอบความใช้ได้ของวิธีและการทดสอบภาคสนามจริง}

คณะวิจัยได้นำระบบต้นแบบไปติดตั้งทดสอบจริงในแปลงสวนทุเรียน 30 แปลงในจังหวัดจันทบุรี และฟาร์มเลี้ยงกุ้งขาว 12 บ่อในจังหวัดตราด ผลการตรวจสอบความใช้ได้ของวิธี (Method Validation)\index{การตรวจสอบความใช้ได้ของวิธี (Method Validation)}\index{Method Validation} และสมรรถนะของระบบสรุปได้ดังนี้

\subsection{ผลการประหยัดพลังงานและการลดต้นทุนในฟาร์มกุ้ง}

การใช้ระบบควบคุมความเร็วรอบ VFD แบบ 4 ลำดับขั้นร่วมกับหัววัด Optical DO สามารถลดการใช้พลังงานไฟฟ้าของเครื่องตีน้ำลงได้เฉลี่ย \textbf{ร้อยละ 41.2} เมื่อเทียบกับการเปิดมอเตอร์แบบความเร็วคงที่ตลอด 24 ชั่วโมง โดยยังคงรักษาระดับออกซิเจนละลายน้ำในบ่อไม่ให้ต่ำกว่า $4.5\text{ mg/L}$ ได้อย่างต่อเนื่อง

ในส่วนของระบบ Smart Autofeeder ที่มีกล้องตรวจยอยกอาหารกุ้ง สามารถลดปริมาณอาหารสูญเสียลงได้ \textbf{ร้อยละ 19.5} ส่งผลให้อัตราการเปลี่ยนอาหารเป็นเนื้อ (FCR) ดีขึ้นจาก $1.52$ เหลือเพียง $1.28$ และลดการสะสมของแอมโมเนียในน้ำลงอย่างมีนัยสำคัญทางสถิติ ($p < 0.01$)

\subsection{ผลการจำแนกความสุกแก่และการคัดเกรดทุเรียน}

ระบบพหุรูปแบบ (Multi-Modal AI) ที่ผสานข้อมูลภาพสีร่องพูทุเรียนเข้ากับความถี่การสั่นพ้องเสียงเคาะ สามารถจำแนกทุเรียนแก่ (เนื้อแห้ง $\ge 32\%$) ได้อย่างแม่นยำสูงถึง \textbf{ร้อยละ 98.4} โดยไม่มีความผิดพลาดในการส่งออกทุเรียนอ่อนหลุดรอดสู่ตู้คอนเทนเนอร์ ช่วยรักษาชื่อเสียงและมูลค่าการส่งออกทุเรียนไทยในตลาดสากล

\section{เกณฑ์การประเมินผลโครงงานบูรณาการ}

โครงงานบูรณาการ (Capstone Project) มีสัดส่วนคะแนนร้อยละ 50 ของหลักสูตร โดยใช้เกณฑ์การประเมินรูบริก 4 มิติ ได้แก่
\begin{enumerate}[leftmargin=1.5em]
    \item \textbf{ความถูกต้องเชิงวิศวกรรมและฟิสิกส์ (ร้อยละ 15):} การประยุกต์ใช้สมการหลักการแรก การต่อวงจรฮาร์ดแวร์ และการปฏิบัติตามมาตรฐานความปลอดภัย
    \item \textbf{การพัฒนาซอฟต์แวร์และปัญญาประดิษฐ์ฝังตัว (ร้อยละ 15):} โครงสร้างโค้ดที่มีประสิทธิภาพ การจัดการข้อผิดพลาด และความแม่นยำของโมเดล AI
    \item \textbf{ผลสัมฤทธิ์ของการทดสอบภาคสนามและการประหยัดต้นทุน (ร้อยละ 10):} ข้อมูล telemetry จริง ผลการประหยัดพลังงานหรือการเพิ่มผลผลิตที่พิสูจน์ได้เชิงประจักษ์
    \item \textbf{การนำเสนอและเอกสารรายงานเชิงวิชาการ (ร้อยละ 10):} การสื่อสารที่ชัดเจน การจัดทำเอกสารตามมาตรฐาน มรภ.รำไพพรรณี และการตอบข้อซักถามของคณะกรรมการ
\end{enumerate}

\section*{คำถามทบทวนท้ายบทที่ 10}
\addcontentsline{toc}{section}{คำถามทบทวนท้ายบทที่ 10}

\begin{enumerate}
    \item จงอธิบายบทบาทของการเชื่อมต่อแบบ End-to-End Pipeline ในการยกระดับสวนผลไม้และฟาร์มสัตว์น้ำเข้าสู่เกษตรแม่นยำระดับ 4.0
    \item การควบคุมเครื่องตีน้ำด้วย VFD แบบ 4 สเต็ป สามารถประหยัดพลังงานไฟฟ้าได้ร้อยละ 40 ได้อย่างไร โดยที่สัตว์น้ำไม่เกิดสภาวะขาดออกซิเจน
    \item จงวิเคราะห์ประโยชน์ของการผสานข้อมูลภาพถ่าย (Vision) ร่วมกับสัญญาณเสียงเคาะ (Acoustics) ในการตัดสินใจจำแนกความสุกแก่ของทุเรียน
    \item หากท่านต้องออกแบบโครงงาน AIoT สำหรับแก้ปัญหาภัยแล้งในสวนผลไม้ชุมชน ท่านจะบูรณาการเทคโนโลยีใดบ้าง พร้อมระบุตัวชี้วัดความสำเร็จเชิงปริมาณ
\end{enumerate}
\clearpage
